\section{Problem Statement and Primary Contributions}

Automatic classification of political ideology in news text is complicated by three challenges that recur across the literature. First, there is no uncontested ground-truth label: ideological judgments are inherently perspectival, and different rating bodies (AllSides, Media Bias/Fact Check, and others) sometimes disagree on the same outlet. Second, the appropriate label granularity is contested. Source-level classification \citep{baly2020acl, darwish2020aaai, ronnback2025biasoutlets} is convenient but collapses meaningful within-outlet variation, while sentence-level classification \citep{spinde2021mbic, spinde2021babe, menzner2025biasscanner} captures local phrasing but loses the cumulative framing that accrues across an article. Third, supervised models fit on article-ideology pairs are vulnerable to source leakage: when the same outlets appear in both training and evaluation, classifiers learn to identify the outlet from stylistic cues and route predictions through its known partisan label rather than engaging with article content \citep{baly2020we}. This thesis adopts article-level classification with explicit source-leakage controls, the granularity at which downstream applications such as media literacy tools, recommendation diversification, and journalistic auditing are most directly benefited.

The inductive bias we add on top of this is \emph{metalinguistic}. Metalinguistic awareness, the capacity to reflect on language as an object of thought rather than merely to use it for comprehension, is a foundational competence in reading \citep{tunmer1984development, nagy2007metalinguistic, Zipke10092007}, and its footprint is visible in ordinary political reporting. Table~\ref{fig:khalil} illustrates the phenomenon: AP and Fox News cover the same arrest, yet their word choices diverge systematically in ways that encode political stance. These are precisely the cues a human annotator draws on when assessing political orientation, and they are the signal our continued pre-training framework targets by attaching per-article theme and tone prediction heads to a contrastive learning backbone designed to learn the encoded ideological stance.

Large language models are increasingly deployed as drop-in classifiers or surrogate annotators for tasks like ideology detection \citep{wu2023large, horton2023large, li2024frontiers, ma-etal-2025-algorithmic}, and they often achieve competitive aggregate accuracy. This convenience introduces a methodological risk that standard evaluation cannot detect. \citet{gao2025LLMscaution} document that LLM behaviour fails to track the human behaviour distribution in several social-science domains, and \citet{yuan2024llmsovercomeshortcutlearning} show that LLMs remain sensitive to lexical shortcuts even at scale. Two classifiers that agree on labels may route predictions through entirely different evidence, and a classifier that exploits a spurious shortcut will register as successful on the test set, but fail to generalize. Diagnosing this gap requires moving beyond input-output evaluation entirely.

A sizable correlational literature establishes that sentiment, thematic emphasis, and topic-directed stance covary with perceived ideology: \citet{smirnova2017ideology} trace tonal shifts in \textit{New York Times} coverage following major geopolitical events, \citet{bhatia2018topic} show that topic-specific sentiment polarities identify ideology in congressional and news text, and \citet{bestvater2023sentiment} document that overall tone and target-directed stance diverge in politically charged writing. What this correlational work cannot settle is whether the relationship is \emph{causal}, for which topics it holds, and whether human and machine annotators route through the same features. We address those questions by usual causal analysis frameworks. Specifically, we estimate topic and community-level average treatment effects via double machine learning, and by decomposing the sentiment-to-ideology pathway via causal mediation analysis, across five annotation paradigms: human AllSides labels, our fine-tuned RoBERTa classifier, a zero-shot GPT-4o-mini baseline, a fine-tuned GPT-4o-mini, and a zero-shot Llama-3.3-70B. Here, the metalinguistic features serve both as the inductive bias that improves classification performance and as the mediator variable through which we trace the causal pathway from text to ideology.

The main finding with this analysis is that zero-shot LLMs systematically inflate effect magnitudes relative to the human baseline, while fine-tuning compresses those magnitudes toward human scale without altering direction. Classification accuracy does not predict causal fidelity: fine-tuned GPT-4o-mini (F1-Macro 72.48) and fine-tuned RoBERTa (F1-Macro 62.19) differ by more than ten F1 points yet produce nearly identical causal pathways overall. Taken together, these results argue that ideology classifiers cannot be certified as trustworthy on F1 alone, and that the same causal framework can be used to compare, calibrate, and diagnose alternative annotation paradigms, including those in which LLMs are used as surrogates for human annotators.

Concretely, this thesis makes five contributions:

\begin{enumerate}
    \item \textbf{Novel metalinguistic integration} of theme and tone features into article-level political bias classification, demonstrating performance improvements particularly for challenging center-leaning content. Our best model, a continued-pretrained RoBERTa augmented with metalinguistic auxiliary objectives, achieves an F1-Macro of 62.19, a 1.59-point improvement over the previous state-of-the-art POLITICS model \citep{liu2022politics} (60.60),
    \item \textbf{Comprehensive LLM evaluation} of state-of-the-art large language models (including GPT-4o-mini, GPT-5-nano, o3-mini, LLaMA-3.1-8B, and LLaMA-3.3-70B) on political bias classification, examining both zero-shot prompting and fine-tuning approaches,
    \item \textbf{Causal inference analysis} demonstrating that annotation paradigm alone can reverse estimated causal directions, that fine-tuning calibrates zero-shot LLM effect magnitudes toward human scale, and that high classification accuracy does not imply causal fidelity. We establish this through a topic- and community-level analysis of theme--tone--ideology relationships across five annotation paradigms (human, fine-tuned RoBERTa, zero-shot GPT-4o-mini, fine-tuned GPT-4o-mini, and Llama-3.3-70B), and position mediation analysis as a diagnostic framework for auditing annotator trustworthiness,
    \item \textbf{Expanded dataset} (AllSidesV2) extending the original AllSides dataset of \citet{baly2020we} from 27,890 articles to 48,290 articles, a 47\% increase in annotated content,
    \item \textbf{Informal human performance baseline} through a controlled annotation study with four participants.
\end{enumerate}
